\label{chap.3}

This chapter describes the parts shared by both generations of the tester: the browser
environment and its action space (Section~\ref{sec:env}), the oracle that decides what is a
fault (Section~\ref{sec:oracle}), and the product layer that turns an application into a
test run and a report (Section~\ref{sec:product}). Figure~\ref{fig:system} shows how they fit
together. The explorer in the middle is the part that differs between v1
(Chapter~\ref{chap.4}) and v2 (Chapter~\ref{chap.5}).

\begin{figure}[H]
\centering
\begin{tikzpicture}[
  box/.style={draw, rounded corners=2pt, align=center, minimum height=1.0cm,
              text width=3.5cm, font=\small, execute at begin node={\hyphenpenalty=10000}},
  arr/.style={-{Latex[length=2mm]}, thick}]
  \node[box] (in) {Application\\\footnotesize URL, static site, Dockerfile repo, suite app};
  \node[box, right=1.0cm of in] (sb) {Intake and sandbox\\\footnotesize policy gate, container};
  \node[box, right=1.0cm of sb] (env) {Browser\\ environment\\\footnotesize Playwright,\\ action registry};
  \node[box, below=1.0cm of env, fill=gray!12] (ex) {Explorer\\\footnotesize v1: AC-DQN\\ v2: hybrid tester};
  \node[box, left=1.7cm of ex] (or) {Oracle\\\footnotesize triggers, seeded faults, replay, controls, ddmin};
  \node[box, left=1.25cm of or] (rep) {Report\\\footnotesize findings, evidence;\\ HTML, MD, JSON};
  \draw[arr] (in) -- (sb);
  \draw[arr] (sb) -- (env);
  \draw[arr] (env.south) ++(-0.4,0) -- ++(0,-1.0) node[midway, left, font=\footnotesize]{page};
  \draw[arr] (ex.north) ++(0.4,0) -- ++(0,1.0) node[midway, right, font=\footnotesize]{action};
  \draw[arr] (ex) -- node[above, font=\footnotesize]{trajectory} (or);
  \draw[arr] (or) -- node[above, font=\footnotesize]{verified} (rep);
\end{tikzpicture}
\caption{The system around the explorer. Only the explorer differs between v1 and v2.}
\label{fig:system}
\end{figure}

\section{The browser environment}
\label{sec:env}

The environment exposes a live application through Playwright and headless Chromium. Each
\emph{episode} starts in a fresh browser context from the application's entry page, so
cookies and local storage do not leak between episodes. For applications behind a login, a
scripted sign-in macro runs at the start of every episode; it is outside the agent's
decisions but inside the interaction budget. Server state is reset between runs, and for the
database-backed suite applications each browser context gets its own data partition
(Chapter~\ref{chap.6}).

After every action the environment waits for the page to settle. It watches the DOM with a
mutation observer and treats the page as settled after a quiet period, rather than waiting
for the network to go idle, which applications that poll or hold a websocket never
satisfy.

\subsection{Action registry}

At each step the registry scans the page and builds the set of available actions, at most
100. It combines page-independent actions (no-op, scroll, resize the viewport, browser
back, forward and refresh) with actions derived from the page's interactive elements:

\begin{itemize}[noitemsep]
  \item \textbf{click} on links, buttons and toggles, and \textbf{rapid click} (a double
        submit) on buttons;
  \item \textbf{type} into text fields, with one value from each of five categories: a
        valid typical value, the boundary minimum, the boundary maximum, an empty string,
        and a value of the wrong type;
  \item \textbf{select} an option of a dropdown;
  \item optionally, \textbf{custom clickables}: elements made interactive by ARIA roles,
        click handlers, tab index or a pointer cursor, which plain HTML scanning misses.
\end{itemize}

When a page has more candidates than slots, the budget is filled in priority order: fixed
actions, then clicks (needed to progress any flow), then valid typed values (so forms can be
submitted), then selects and rapid clicks, then the boundary and invalid variants. Elements
visible in the viewport are allocated before those below the fold. Actions are identified by
stable selectors (an element's id or name, else a structural path), so the same control
keeps its identity across visits.

\subsection{Deterministic triggers}

On every step the environment records generic failure signals: console errors, uncaught
page errors, unexpected HTTP failures (4xx and 5xx responses to the application's own
requests), a page that never settles, and broken navigation (a navigational click that
leaves the URL unchanged). A link back to the page the agent is already on is not
navigational, so it does not count as broken. These triggers are the implicit oracle that
published RL web testers use. Here they feed the reward and the crash-type faults; on their
own they are never reported as verified findings.

\section{The verified-fault oracle}
\label{sec:oracle}

\subsection{Seeded faults and assertions}

Every evaluation application is written with a list of \emph{seeded faults}. Each fault is
a literal source patch: an anchor string that must occur exactly once in a file, and its
replacement. Applying every patch to the clean source gives the \emph{faulty} build that
agents test. For each fault $F$ there is also a \emph{control build}: the faulty build
without $F$'s patch.

Each fault declares how it shows itself, as an assertion over the recorded trajectory. The
assertion language is deliberately closed, with eight kinds:

\begin{itemize}[noitemsep]
  \item \code{text\_equals}, \code{text\_contains}, \code{text\_absent}: a named page region
        must equal, contain or not contain a literal;
  \item \code{echoes\_input}: a value the agent typed earlier must appear on a later page
        (the core check for corrupted state);
  \item \code{numeric\_equals}: a number on the page must equal an arithmetic expression
        over recorded inputs (for example, a total over quantities);
  \item \code{url\_matches}: the step's URL must match a pattern (for broken transitions);
  \item \code{reached\_with\_invalid\_input}: a guarded state was reached although the input
        that guards it was invalid;
  \item \code{region\_empty}: a region the application's contract says must carry a value is
        empty.
\end{itemize}

Regions are named in the application's manifest, so an assertion never contains a CSS
selector. Assertions are evaluated on the recorded pages without a browser, so checking is
a pure, reproducible function of the record. The development suite added four extension
kinds: a region stable across reload, an input absent after submission, arithmetic between
regions, and \code{console\_or\_request\_error} for crash-type faults (Section~\ref{sec:crash}).

\subsection{From candidate to verified finding}

A violation of a fault's assertion during exploration is only a \emph{candidate}.
Verification has three steps:

\begin{enumerate}
  \item \textbf{Replay.} The trajectory is replayed twice in fresh browser sessions against
        the faulty build. The violation must recur both times.
  \item \textbf{Per-fault control.} The same trajectory is replayed against the control
        build without $F$. If the violation still occurs, it is not caused by $F$ (another
        fault or the application's normal behaviour produces it), and the candidate is
        rejected.
  \item \textbf{Minimisation.} Delta debugging~\cite{zeller2002simplifying} removes steps
        while the violation still reproduces, giving a short reproduction.
\end{enumerate}

Only candidates that pass all three are \emph{verified findings}. The rule makes false
positives measurable: a run on the \emph{clean} build has no seeded faults, so any verified
finding there would be a false positive. Every evaluation in this thesis includes such
clean-build control runs, and none produced a verified finding. Zero observed false
positives is not a zero false-positive rate. The held-out evaluation, for example, had 32
clean-build controls; with none positive, a one-sided 95\% upper bound on the per-run rate
is about $3/32 \approx 9\%$ (the rule of three).

\subsection{What the oracle solves, and what it does not}
\label{sec:oracle-limits}

The \emph{oracle problem} is deciding whether observed behaviour is
correct~\cite{barr2015oracle}. For the seeded faults, Katalyze does not solve it. Each
seeded fault comes with an assertion written by whoever seeded it, and verification
establishes that a failure of that assertion is reproducible and caused by that fault. The
oracle makes reported findings \emph{credible}. It does not let the tester tell right from
wrong in an application it has no specification for. Three consequences follow.
\begin{enumerate}[noitemsep]
  \item The evaluations measure how well a tester drives an application into a state where
        a known assertion fails. They do not measure whether it would recognise a bug that
        nobody wrote an assertion for.
  \item On applications without seeded faults, which is what the product mostly sees
        (Section~\ref{sec:product}), findings can come only from generic triggers and the
        specification-free invariants below. Generic triggers are script errors, failed
        requests, error pages and navigation that does nothing. Both kinds are weaker than
        a seeded assertion, and the invariants raised false alarms (Section~\ref{sec:heldout}).
  \item A per-fault control build exists only because the fault was seeded. For a real
        bug there is no control, and verification reduces to replay and minimisation.
\end{enumerate}

\subsection{Crash-type faults}
\label{sec:crash}

The external benchmark (Chapter~\ref{chap.7}) counts failures as console and request
errors, which seeded assertions over page content do not capture. A subset of the
development applications therefore also seeds \emph{crash-type faults}: patches that make a
specific interaction raise an uncaught script error or a failing request. Their assertion,
\code{console\_or\_request\_error}, matches the error's signature. Verification is the
same: the error must recur on replay and be absent from that fault's control build, so a
generic console error elsewhere in the application is not attributed to the fault.

\subsection{Specification-free invariants}

Seeded assertions verify only what was seeded. For arbitrary applications, the v2 tester
also checks three invariants that need no specification: a value the user entered must
survive a reload; it must not come back altered; and the same field on the same page must
not be rejected once and accepted another time. Their candidates are replayed twice and
attributed to seeded faults through the control builds. They are reported separately and
never counted as verified seeded faults.

\section{The product layer}
\label{sec:product}

\paragraph{Intake and restricted execution sandbox.} An uploaded repository must contain a Dockerfile. Its build
definition is parsed and passed through a security policy gate before anything runs. The
gate refuses privileged containers, host mounts, host networking and Compose features that
would merge in a second, unvalidated file. The container then runs with all capabilities
dropped, a read-only root filesystem, no new privileges, CPU, memory and process limits,
and an internal-only network. This is a \emph{restricted} execution sandbox: it narrows what
uploaded code can do, but it has not been tested against deliberately malicious code. The
product is meant for code its user would be willing to run locally.

\paragraph{Test sessions and reports.} A test session runs the chosen agent within an
interaction budget, verifies candidates as above, and writes a report with three parts:
verified findings with minimal reproductions; unverified candidates; and, separately, any
observations from the language model, marked unverified. Page bodies and frames are stored
content-addressed as the run proceeds, so every finding can be traced to what produced it.
Reports are available in HTML, Markdown and JSON.

\paragraph{Interfaces.} A command-line tool tests a URL, a static site, a Dockerfile
repository or a bundled application. A web interface (FastAPI with server-rendered
templates) offers the same as a guided flow. The user picks a demo application, enters a
URL or uploads an archive; chooses an agent (the RL agents or the random baseline) and a
budget; watches a live event stream of the run; and opens or downloads the report. A
research page shows every evaluation and development round, read live from the result
files. The interface binds to the local machine by default, because it builds and runs
uploaded code.

\paragraph{End-to-end acceptance.} Table~\ref{tab:acceptance} lists one end-to-end run of
the final product for each kind of target. Each run was started through the web
interface's own API, ran \code{hybrid\_final}, and produced HTML, Markdown and JSON reports.
For the bundled applications the seeded faults are known, so their findings can be checked.
For the other targets no ground truth exists, and the findings are reproduced and
minimised observations, not confirmed bugs.

\begin{table}[H]
\centering
\small
\caption{End-to-end acceptance runs of the final product (24 October 2026). ``Findings'' are
the verified findings in the report. For targets without seeded faults they come from
generic triggers.}
\label{tab:acceptance}
\footnotesize
\begin{tabularx}{\textwidth}{@{}>{\raggedright\arraybackslash}p{2.6cm} >{\raggedright\arraybackslash}p{2.5cm} r r >{\raggedright\arraybackslash}X@{}}
\toprule
Target kind & Example & Budget & Time & Findings \\
\midrule
Bundled app, full stack in Docker & Helpline & 600 & 7 min & 4: seeded HELP-2, HELP-3 and HELP-4, and one reproduced ``Move down'' control that does nothing \\
Bundled app, full stack in Docker & Stackroom & 600 & 5 min & 2: seeded STK-3, and one reproduced script error \\
Bundled app, full stack in Docker & Orbit & 600 & 5 min & 1: seeded ORB-3 \\
Bundled benchmark app & Fernbrook Outfitters & 1{,}000 & 4 min & 1: seeded SHOP-01 \\
Static site, bundled demo & Harbor \& Pine & 300 & 5 min & 4 (curiosity arm): missing pages, a broken form, a script error \\
Static site, zip upload & Harbor \& Pine & 200 & 3 min & 4: the same kinds \\
Dockerfile repository & Northwind static site & 200 & 3 min & 2: a missing page and a script error \\
Public URL & TodoMVC (demo.\-playwright.dev) & 200 & 3 min & 0 (before the fix below: 3 false reports on links to other sites) \\
\bottomrule
\end{tabularx}
\end{table}

The acceptance runs found three product defects, all now fixed. The first was fixed before
every run in the table. The TodoMVC run is the only one made after the other two; no other
run in the table reported a finding on a link to another site.
\begin{itemize}[noitemsep]
  \item A form that submits and returns to the same address with updated content (Helpline's
        post-redirect-get forms) was reported as navigation that does nothing. Such a report
        now requires the page text to stay the same as well as the address.
  \item Links to other websites, which the tester refuses to follow, were reported the same
        way. They are now excluded.
  \item One slow page load at the start of an episode ended the whole session on a public
        site. The opening load is now retried up to three times.
\end{itemize}
